\section{Effective synthesis and uniform compilation}
\label{syn:section}
\label{sec:synthesis}

This appendix gives a general route from a finite tensor saving to a
uniform Fourier algorithm.  Its purpose is to separate two questions:
how to obtain usable coefficients from an existence theorem over $\C$,
and how to use effective small circuits when their compilation time has
no useful bound.  The explicit network in the main argument gives the
stronger quantitative conclusion without this search.
We first replace a complex certificate by identities in a fixed
finite-dimensional rational algebra.  Matrix representations of that
algebra then give ordinary scalar circuits without choosing a complex
embedding.  We print those circuits together with their scalar
preparation programs, and finally limit the printer's running time to
select a sufficiently small compiled radix for each input length.

The certificate used below consists of two finite words for the same
invertible nonmonomial matrix $A\in\GL_q(\C)$, with integers
$q,b\ge2$. The first word acts on exactly
$q^b$ coordinates as $A^{\otimes b}$ and uses
\begin{equation}\label{syn:win}
 g<bq^{b-1}
\end{equation}
forward calls to $A$, interspersed with invertible monomials. A call
acts on an ordered $q$-tuple and fixes every other coordinate. The
second word acts on $q$ coordinates as $U=I_q+E_{12}$, where $E_{12}$
is a matrix unit, using a positive number of forward calls to that same
$A$ and invertible monomials.

The companion paper~\cite[Theorem~8.3 and Lemma~2.2]{exactfourier}
obtains such a pair from a general finite-win theorem and positive
generation by an invertible nonmonomial matrix. The following lemma
supplies the particular certificate needed here directly from our
network. Thus the search-termination proof is also available locally.
The extraction and printer constructions that follow use only the two
word identities, so they retain their general certificate interface.

\begin{lemma}[A local certificate for the symbolic search]
\label{syn:local-certificate}
For the invertible nonmonomial matrix $C$ in \eqref{net:C}, there are
integers $b\ge2$ and $g\ge0$ and a word on exactly $2^b$ coordinates
which implements $C^{\otimes b}$ using $g<b2^{b-1}$ forward calls to
$C$ and invertible monomials. On two coordinates, $I_2+E_{12}$ has a
word with exactly three forward calls to the same $C$ and invertible
diagonal factors.
\end{lemma}

\begin{proof}
Use Proposition~\ref{net:finite-interface} and the padding
\eqref{net:padding} on $W_*=2^{71}$ arrays of length $2^{mf}$.
The $S=W_*m-\Delta$ directional steps require exactly
$Sf2^{mf-1}$ forward $C$-calls when their tensor factors are
applied separately: one step has $2^{(m-1)f}$ fibers, each requiring
$f2^{f-1}$ calls. Inverse steps add only translations.
Each scalar row of \eqref{net:schedule} is a composition of elementary
shears between existing, distinct roles; its source coordinates are
unchanged within that row. The coefficients are nonzero rational
numbers. Equation~\eqref{loc:nonzero-shear} therefore bounds the
pointwise work by $H2^{mf}$ forward calls for a fixed integer $H$
independent of $f$. All remaining operations are invertible monomials.
Every role carries arbitrary data, so this word uses exactly
$W_*2^{mf}$ coordinates and acts as
$I_{W_*}\otimes C^{\otimes mf}$. Address permutations here are
monomial factors on these coordinates; no implementation buffer is
added to the word's width.

Identify the role index with 71 binary axes and apply $C$ along
each of them. This costs $71W_*2^{mf-1}$ calls and gives
$C^{\otimes b}$ with $b=mf+71$. Its actual number $g$ of forward calls
therefore satisfies
\[
 g\le (Sf+71W_*+2H)2^{mf-1}
   =b2^{b-1}-(\Delta f-2H)2^{mf-1}<b2^{b-1}
\]
for any integer $f>2H/\Delta$.
Finally, $\det C=i$ and both scalar coefficients in \eqref{net:C} are
nonzero, so $C$ is invertible and nonmonomial.
Equation~\eqref{loc:nonzero-shear} with $t=1$ supplies the claimed
three-call word for $I_2+E_{12}$. Any embedded call on an ordered pair
is a permutation conjugate of $C\oplus I_{2^b-2}$; the permutations
can be absorbed into adjacent monomial factors.
\end{proof}

\subsection{Obtaining finite rational algebra data}

The extraction uses the triangular monic substitution from the standard
proof of Noether normalization~\cite[Tags 051M, 051N, and 00OX]{stacksNormalization}
and the determinant form of Nakayama's lemma~\cite{stacksNakayama}.
We give the rational algorithm and its termination proof in full.

\begin{lemma}[Effective finite algebra extraction]\label{syn:extraction}
Given finitely many polynomials $f_i$ in $a$ variables over $\Q$, a
deterministic rational algorithm either certifies that $(f_i)$ is the
unit ideal or produces a nonzero finite-dimensional commutative unital
$\Q$-algebra $E$, with a rational multiplication table, and elements
$\xi_1,\ldots,\xi_a\in E$ satisfying every $f_i(\xi)=0$.
In particular, a complex solution guarantees the second outcome.
Required units can be imposed by adding inverse-witness equations.
\end{lemma}

\begin{proof}
Keep a finite generating list for the ideal $I=(f_i)$ and an integer $m$,
initially $m=a$.  Retain polynomial-combination witnesses for derived
members of $I$; rational elimination and the monic reductions below
can track these witnesses.
After invertible polynomial changes of coordinates, maintain, for
each $j>m$, a known member of $I$ of the form
\[
 x_j^{d_j}+\sum_{e<d_j}c_{je}(x_1,\ldots,x_{j-1})x_j^e,
 \qquad d_j\ge1.
\]
Thus $S=\Q[x_1,\ldots,x_a]/I$ is finite as a module over
$R=\Q[x_1,\ldots,x_m]$.

First compute the fiber $S/(x_1,\ldots,x_m)S$.  This is effective:
set $x_1=\cdots=x_m=0$ and let $B$ be the quotient by only the
displayed monic relations.  As a tower of monic extensions, $B$ has
a rational basis consisting of the monomials whose $x_j$-exponents
are less than $d_j$ for $j>m$.  Successive monic reduction gives its
multiplication table.  If $\bar f_i$ are the images in $B$ of the
generators of $I$, form the rational span of all $b\bar f_i$ with
$b$ in this basis.  Since the basis spans $B$, this is exactly the
ideal $\sum_i B\bar f_i$.  Its quotient is the required fiber, so
rational linear algebra computes the fiber and its table.  A nonzero
fiber is the requested answer, after undoing the coordinate changes.
If $m=0$ and the fiber is zero, then $S=0$, so $I$ is the unit ideal.

Suppose the fiber is zero and $m>0$.  We claim $I\cap R\ne0$.
Otherwise $R\hookrightarrow S$.  Localize at
$\mathfrak m=(x_1,\ldots,x_m)$ in $R$.  The resulting finite
$R_{\mathfrak m}$-module $S_{\mathfrak m}$ is nonzero, since it still
contains $R_{\mathfrak m}$.  The zero fiber would give
$S_{\mathfrak m}=\mathfrak mS_{\mathfrak m}$.  For a finite generating
column $v$ this means $v=Mv$ with all entries of $M$ in the maximal
ideal.  Multiplication by the adjugate of $I-M$ gives
$\det(I-M)v=0$; its determinant is a unit, a contradiction.

Find a nonzero member $h\in I\cap R$ by enumerating increasing degree
cutoffs for monomial multiples of the generators of $I$.  At each
cutoff, rational linear algebra tests whether their span has a
nonzero element involving only the first $m$ variables.  This search
terminates by the claim and the definition of an ideal.  A nonzero
constant certifies $I=(1)$.  Otherwise let $d=\deg h\ge1$ and substitute
\[
 x_i=y_i+y_m^{(d+1)^i}\quad(1\le i<m),\qquad x_m=y_m,
\]
leaving later variables unchanged.  This change is invertible.
The weights $1,(d+1),\ldots,(d+1)^{m-1}$ distinguish all monomials of
total degree at most $d$ by base-$(d+1)$ expansion.  Consequently the
highest power of $y_m$ in the transformed $h$ has a nonzero rational
leading coefficient: it comes from a unique monomial, taking the
pure-power term in every substituted factor.  Divide by that
coefficient to obtain a monic relation in $y_m$.  Earlier triangular
relations remain monic in their respective later variables.  Rename
the variables and decrease $m$ by one.  At most $a$ such decreases
occur, proving termination.  Every step uses finite rational
arithmetic.  Finally, the equation $uv=1$ forces $u$ to be a unit in
every output algebra, including one with zero divisors.
\end{proof}

Enumerate integer tuples $(q,b,g,h)$ with $q,b\ge2$, $g\ge0$, $h\ge1$ and
\eqref{syn:win}.  For each tuple form polynomial equations for
\begin{equation}\label{syn:certificate}
 A^{\otimes b}=M_0YM_1Y\cdots YM_g,
 \qquad U=N_0AN_1A\cdots AN_h,
 \qquad Y=A\oplus I_{q^b-q}.
\end{equation}
All embedded placements are absorbed into adjacent monomials.
Include a matrix inverse witness for $A$.  In a monomial slot write
a variable permutation factor $P$ and a variable diagonal factor;
impose $p_{ij}^2=p_{ij}$ and row and column sums equal to $1$ on $P$,
and give every diagonal entry an inverse witness.  These are finite
polynomial systems over $\Q$.  Apply Lemma~\ref{syn:extraction} in a
fixed enumeration order until one succeeds.  Lemma~\ref{syn:local-certificate} gives a complex solution for a tuple
with $q=2$ and $h=3$, so success is guaranteed.

The entries $p_{ij}$ in the resulting $E$ may be nonscalar idempotents.
Resolve them one at a time.  For a current idempotent $p\ne0$, replace
$E$ by the nonzero algebra $pE$ with unit $p$ and project every stored
element by $z\mapsto pz$; if $p=0$, do nothing.  This is a unital
homomorphism onto the new algebra and makes the current idempotent
equal to its unit.  Previously resolved values remain $0$ or $1$.
After finitely many rational table operations all permutation entries
are $0$ or $1$, and their row and column sums give actual permutations.
All identities, inverse witnesses, and diagonal units survive.
Choose a rational basis of the final nonzero algebra beginning with
$1_E$.  We have thus computed fixed, fully specified data
\eqref{syn:certificate} over $E$.

\begin{remark}[Faithful scalar simulation]\label{syn:simulation}
If $d=\dim_{\Q}E$, then $E\otimes_{\Q}\C$ is the vector space $\C^d$
with the computed multiplication table.  The embedding
$z\mapsto z1_E$ is injective.  Addition is coordinatewise and
multiplication by any prepared algebra element is a $d\times d$
complex matrix, so each linear algebra operation costs $O(d^2)$
ordinary scalar operations.  When a circuit acts as $F_r$ times the
algebra identity, embed its inputs as multiples of $1_E$ and read the
first coordinate of each output.  This is a faithful, constant-factor
simulation; choosing a complex embedding or extracting algebraic
roots is unnecessary.  Here $d$ is fixed before the input length is
known.
\end{remark}

\subsection{A printer for increasingly efficient radices}

We record why the fixed identities suffice for an effective circuit
printer.  The tensor amplification of \cite[Lemma~2.3]{exactfourier}
also works over $E$: its proof uses only the identities and linear
operations.  Explicitly, put $Q=q^b$, $j=\lfloor k/b\rfloor$ and
tensor the first word in \eqref{syn:certificate} across $j$ groups.
Every call slot decomposes, after reindexing, as
\[
 (A\oplus I_{Q-q})^{\otimes j}
 \cong\bigoplus_{r=0}^j
 (A^{\otimes r})^{\oplus\binom jr(Q-q)^{j-r}}.
\]
Recurse on these smaller tensor orders and use direct calls on the
fewer than $b$ remaining axes.  Monomial scalings and copies cost
linear size at each level.  Let $f(k)$ be the size of this recursive
construction divided by $q^k$.  The total width of the sectors with
$r$ active factors, divided by $Q^j$, is
\[
 \frac{\binom jr(Q-q)^{j-r}q^r}{Q^j}
 =\binom jr(q/Q)^r(1-q/Q)^{j-r}.
\]
Thus the normalized cost satisfies
\[
 f(k)\le C_0+g\,\mathbb E f(K),\qquad
 K\sim\operatorname{Bin}(\lfloor k/b\rfloor,q/Q).
\]
The distribution weights sectors by their dimensions.  Since
$gq/(bQ)<1$, choose $0<\alpha<1$ with
$g(q/(bQ))^\alpha<1$.  Concavity and induction, using
$\mathbb E(K+1)^\alpha\le(1+kq/(bQ))^\alpha$, give
\begin{equation}\label{syn:amplification}
 f(k)=O((k+1)^\alpha).
\end{equation}
The strict inequality absorbs $C_0$ for large $k$; enlarge the
inductive constant for the finite initial range.  Neither $\alpha$
nor that constant is needed by the printer.

For successive $t\ge1$, let $r_1<\cdots<r_t$ be the first $t$ primes
greater than $2q$, and put
\[
 m_t=\prod_{i=1}^t r_i,\qquad
 \ell_t=m_t\,2^{\lceil\log_2(16r_t)\rceil},\qquad
 R_t=\Q[z]/(\Phi_{\ell_t}(z)).
\]
Cyclotomic polynomials and arithmetic in this finite algebra are
effective: compute $\Phi_s$ recursively from
$z^s-1=\prod_{d\mid s}\Phi_d(z)$, and use polynomial division and
the Euclidean algorithm.  A polynomial is a unit modulo $\Phi_s$
exactly when their gcd is $1$.  The distinct roots of $\Phi_s$ are
the primitive $s$th roots, so identities in $R_t$ can be checked at
all primitive specializations.  Roots of orders dividing $\ell_t$
are compatible powers of $z$.
Carry out the next constructions over $E\otimes_{\Q}R_t$, with the
Fourier coefficients in the central factor $R_t$.  The two certificate
identities and their inverse witnesses persist under extension of
scalars; in particular, every certified unit remains a unit.

Use the exact-width Fourier layer construction of Lemma~\ref{loc:fourier}
on each width $r_i$.  It gives $O(1+\log^4 r_t)$ monomial or disjoint
shear layers.  Its rational formulas hold in $R_t$: all denominators
are nonzero at every primitive specialization, and the convolution
root orders divide the indicated power of two.  In particular the
Newton denominators $\prod_{a=1}^j(1-u^a)$ have $j<r_i$ for a
primitive $r_i$th root $u$.  The measured workspace test is a test on
the integer circuit description and is unchanged by this scalar
extension.

To substitute the second word of \eqref{syn:certificate}, split any
shear coefficient $c\in R_t$ into two units $d+(c-d)$, where $d$ is
a positive integer chosen by testing $d=1,2,\ldots$ until $c-d$ is a
unit.  This terminates: only finitely many complex values of $c$ at
the primitive roots can be positive integers.  A shear with a unit
coefficient $v$ is obtained by conjugating $U$ by
$\diag(v,1,\ldots,1)$, so it has the fixed word pattern from
\eqref{syn:certificate}.  The resulting word restores the arbitrary values of
its $q-2$ spectator coordinates.  Partition the disjoint pairs of an
axis into at most $q$ batches of at most $\lfloor r_i/q\rfloor$ pairs;
this is possible because $r_i>2q$ gives
$q\lfloor r_i/q\rfloor\ge\lfloor r_i/2\rfloor$.
Each batch fits disjoint $q$-tuples inside the same $r_i$ coordinates.
Its spectators may belong to other pairs, since they are restored
before the next batch.

Synchronize these fixed word patterns across axes, with identities
in idle positions.  Their product is the tensor of the local products
by the tensor multiplication identity.  Each monomial slot has linear
size.  Every other slot is a direct sum of $A^{\otimes k}$ with
$k\le t$ and total width $m_t$, so \eqref{syn:amplification} bounds
its size by $O(m_t(t+1)^\alpha)$.  Chinese remainder reindexing gives
$F_{m_t}$ from the tensor of the local Fourier transforms.
Now specialize $R_t$ by $z\mapsto\zeta_{\ell_t}$, obtaining a circuit
over $E\otimes_{\Q}\C$ whose output matrix is $F_{m_t}$ times the
algebra identity.  Expand only the fixed algebra $E$ by
Remark~\ref{syn:simulation}.  Its dimension contributes a constant
factor, whereas $R_t$ has supplied polynomial expressions for the
complex coefficients.  The resulting complex linear circuit has size
\begin{equation}\label{syn:radix-size}
 s_t=O\bigl(m_t(t+1)^\alpha(1+\log^4(t+2))\bigr)
     =o(m_t\log m_t).
\end{equation}
Here $r_t=O_q((t+1)^2)$, by the elementary prime estimate in
Lemma~\ref{lem:prime-lengths}, and $\log m_t\ge t\log2$.  The printer computes all
tensor coefficients and all wiring explicitly.  It may take any
finite amount of bit computation to do so.  Each coefficient is a
polynomial with rational coefficients evaluated at
$\zeta_{\ell_t}$; symbolic inverses have already been reduced in
$R_t$.  For each polynomial coefficient, print a finite Horner evaluation list,
with its rational literals written explicitly; sharing previously computed
values is optional.  Thus the printed record contains a complete circuit
and a fully expanded scalar preparation program, with no unspecified
complex constants or hidden loops.

\subsection{From a slow printer to a fast uniform algorithm}

\begin{theorem}[Bounded-printer uniformization]\label{syn:printer}
Suppose a fixed bit machine prints a sequence of complete records
$(m_t,C_t)$ with strictly increasing $m_t\to\infty$, where $C_t$
computes $F_{m_t}$ by $s_t=o(m_t\log m_t)$ scalar linear gates.
Each record must explicitly list the gates, output indices, and a
scalar preparation program with every instruction explicitly listed
and executed once in order.  Its instructions use rational literals, field
operations with valid nonzero divisors, and canonical Fourier roots
$\zeta_a=\exp(2\pi i/a)$ of positive integer orders
whose orders are written literally in binary.  Then there is one
deterministic algorithm computing every $F_n$ in $o(n\log n)$ time,
including preparation and indexing in the model of this paper.
A single specified root of order less than $8n^3$ suffices.
\end{theorem}

\begin{proof}
The printing budget will control both the circuit description and all
literal scalar data, without requiring a bound on when a record appears.
Normalize the printer to output at most one bit per machine step.
For $n\ge2$, run it from the beginning for
$B=\lfloor\log_2 n\rfloor$ bit steps.  Select the last completed
record with $2\le m_t\le\sqrt B$; if there is none, use the ordinary radix
$m=2$ circuit.  Every fixed record eventually finishes and satisfies
the width bound, so the selected stream index tends to infinity.
In particular $m\to\infty$ and $s/(m\log m)\to0$.

The total visible description has at most $B$ bits.  Decode the
selected explicit lists to addressable arrays and evaluate its
preparation program once.  Construct rational literals by their
binary expansions and perform each listed operation in the stated
arithmetic model.  This costs $B^{O(1)}$, including parsing and bit
simulation.  The gate list then executes on each new input in
$O(m+s)$ time, including entry and output copies.  There is no
expansion of a succinct circuit during its invocation.

Put $X=2n-1$, let $m^H$ be the largest power of $m$ not exceeding
$X$, and set
\[
 p=\lceil X/m^H\rceil,\qquad P=m^Hp.
\]
Then $1\le p\le m$, $X\le P<2X<4n$, and
$H\le\log(2n)/\log m$.  The mixed-radix Fourier recursion uses $H$
copies of radix $m$ and one radix $p$, the latter computed directly.
For clarity, at a split $ab$ use input index $bj_a+j_b$ and output
index $k_a+ak_b$.  The identity
\[
 (bj_a+j_b)(k_a+ak_b)
 \equiv bj_ak_a+aj_bk_b+j_bk_a\pmod {ab}
\]
gives the two smaller Fourier transforms and their intervening
twiddle.  At each level, copying subarrays into the required order,
tabulating or reading twiddles, and enumerating the two indices cost
linear work.  Reuse the decoded radix circuit each time.  Tabulate
powers of $\zeta_P$ in $O(P)$ operations; every recursive width
divides $P$.  Thus one transform, or its inverse by frequency
reversal and scaling, costs
\begin{equation}\label{syn:uniform-cost}
 O\bigl(P[1+p+H(1+s/m)]\bigr)+B^{O(1)}.
\end{equation}
All indices and products of indices have polynomial bounds in $n$,
so a constant number of $O(\log(n+2))$-bit words suffices per access.

Finally reduce $F_n$ to convolution of length $P$.  With
$\eta=\zeta_{2n}$, store $\eta^{j^2}x_j$ for $0\le j<n$ and pad
with zeros.  The fixed kernel stores $\eta^{-h^2}$ at the residues
of $-(n-1)\le h\le n-1$; these residues are distinct since
$P\ge2n-1$.  Multiply convolution output $k$ by $\eta^{k^2}$.
The identity $k^2+j^2-(k-j)^2=2kj$ gives exactly $F_nx$.
Computing the fixed kernel's transform, the variable transform and
the inverse uses only a constant number of transforms bounded by
\eqref{syn:uniform-cost}, plus linear preparation and scaling work.
The chirp scalars come from an $O(n)$-size table of powers of $\eta$,
indexed by the squared exponents modulo $2n$.

For the root accounting, let $R$ be the product of the root orders
listed in the selected preparation program, or $1$ if there are
none.  Their total binary length is at most $B$, whence
$R\le2^B\le n$.  The single root
$\zeta_D$, where $D=2nPR<8n^3$, supplies every listed root and the
two outer roots by integer powering.  There are at most $B$ such
requests, and powering and forming their exponents cost $B^{O(1)}$.
This includes obtaining the required powers of the supplied root in the
preparation bound.

Since $p=O(1+\sqrt{\log n})$, dividing the arithmetic term of
\eqref{syn:uniform-cost} by $n\log n$ gives
\[
 O\left(\frac{1+p}{\log n}
       +\frac1{\log m}+\frac{s}{m\log m}\right)=o(1).
\]
The polynomial logarithmic preparation term is also $o(n\log n)$.
Length $1$ is the identity.  This proves the claim, with all choices
independent of the input values.
\end{proof}

Applying Theorem~\ref{syn:printer} to the printer above yields a
uniform all-length $o(n\log n)$ algorithm directly from the finite
certificate interface.  The symbolic search, finite algebra, and stored
circuits are implementation devices; the only supplied complex
primitive is the explicitly specified Fourier root.  This conclusion
complements the explicit quantitative construction of the main text.
